\documentclass[10pt,a4paper]{amsart}
\usepackage{amsmath,amssymb,amsthm,graphicx}
\usepackage{verbatim}
\newenvironment{reference}{\comment}{\endcomment}
%\newenvironment{reference}{}{}
\newenvironment{slogan}{\comment}{\endcomment}
\newenvironment{history}{\comment}{\endcomment}

% For commutative diagrams we use Xy-pic
\usepackage[all]{xy}

% We use 2cell for 2-commutative diagrams.
\xyoption{2cell}
% Theorem environments.
%
\theoremstyle{plain}
\newtheorem{theorem}[subsection]{Theorem}
\newtheorem{proposition}[subsection]{Proposition}
\newtheorem{lemma}[subsection]{Lemma}

\theoremstyle{definition}
\newtheorem{definition}[subsection]{Definition}
\newtheorem{example}[subsection]{Example}
\newtheorem{exercise}[subsection]{Exercise}
\newtheorem{situation}[subsection]{Situation}

\theoremstyle{remark}
\newtheorem{remark}[subsection]{Remark}
\newtheorem{remarks}[subsection]{Remarks}

\numberwithin{equation}{subsection}

% Macros
%
\def\lim{\mathop{\mathrm{lim}}\nolimits}
\def\colim{\mathop{\mathrm{colim}}\nolimits}
\def\Spec{\mathop{\mathrm{Spec}}}
\def\Hom{\mathop{\mathrm{Hom}}\nolimits}
\def\Ext{\mathop{\mathrm{Ext}}\nolimits}
\def\SheafHom{\mathop{\mathcal{H}\!\mathit{om}}\nolimits}
\def\SheafExt{\mathop{\mathcal{E}\!\mathit{xt}}\nolimits}
\def\Sch{\mathit{Sch}}
\def\Mor{\mathop{\mathrm{Mor}}\nolimits}
\def\Ob{\mathop{\mathrm{Ob}}\nolimits}
\def\Sh{\mathop{\mathit{Sh}}\nolimits}
\def\NL{\mathop{N\!L}\nolimits}
\def\CH{\mathop{\mathrm{CH}}\nolimits}
\def\proetale{{pro\text{-}\acute{e}tale}}
\def\etale{{\acute{e}tale}}
\def\QCoh{\mathit{QCoh}}
\def\Ker{\mathop{\mathrm{Ker}}}
\def\Im{\mathop{\mathrm{Im}}}
\def\Coker{\mathop{\mathrm{Coker}}}
\def\Coim{\mathop{\mathrm{Coim}}}

% Boxtimes
%
\DeclareMathSymbol{\boxtimes}{\mathbin}{AMSa}{"02}

%
% Macros for moduli stacks/spaces
%
\def\QCohstack{\mathcal{QC}\!\mathit{oh}}
\def\Cohstack{\mathcal{C}\!\mathit{oh}}
\def\Spacesstack{\mathcal{S}\!\mathit{paces}}
\def\Quotfunctor{\mathrm{Quot}}
\def\Hilbfunctor{\mathrm{Hilb}}
\def\Curvesstack{\mathcal{C}\!\mathit{urves}}
\def\Polarizedstack{\mathcal{P}\!\mathit{olarized}}
\def\Complexesstack{\mathcal{C}\!\mathit{omplexes}}
% \Pic is the operator that assigns to X its picard group, usage \Pic(X)
% \Picardstack_{X/B} denotes the Picard stack of X over B
% \Picardfunctor_{X/B} denotes the Picard functor of X over B
\def\Pic{\mathop{\mathrm{Pic}}\nolimits}
\def\Picardstack{\mathcal{P}\!\mathit{ic}}
\def\Picardfunctor{\mathrm{Pic}}
\def\Deformationcategory{\mathcal{D}\!\mathit{ef}}
\usepackage{cleveref}
\crefname{originallemma}{original lemma}{original lemmas}
\Crefname{originallemma}{Original lemma}{Original lemmas}
\makeatletter
\newlabel{lemma-restriction-flat}{{0E8J}{source}{}{sourcepointer-1}{}}
\newlabel{lemma-restriction-flat@cref}{{[originallemma][0][]0E8J}{[1][0][]source}{}{}{}}
\makeatother
\crefname{originalreference}{original reference}{original references}
\Crefname{originalreference}{Original reference}{Original references}
\makeatletter
\newlabel{algebra-lemma-relations}{{04VX}{source}{}{sourcepointer-2}{}}
\newlabel{algebra-lemma-relations@cref}{{[originalreference][0][]04VX}{[1][0][]source}{}{}{}}
\makeatother
\crefname{originalreference}{original reference}{original references}
\Crefname{originalreference}{Original reference}{Original references}
\makeatletter
\newlabel{algebra-lemma-module-colimit-fp}{{00HA}{source}{}{sourcepointer-3}{}}
\newlabel{algebra-lemma-module-colimit-fp@cref}{{[originalreference][0][]00HA}{[1][0][]source}{}{}{}}
\makeatother
\crefname{originalreference}{original reference}{original references}
\Crefname{originalreference}{Original reference}{Original references}
\makeatletter
\newlabel{algebra-lemma-tensor-products-commute-with-limits}{{00DD}{source}{}{sourcepointer-4}{}}
\newlabel{algebra-lemma-tensor-products-commute-with-limits@cref}{{[originalreference][0][]00DD}{[1][0][]source}{}{}{}}
\makeatother
\title{Equational criterion of flatness for modules on a ringed site}
\author{The Stacks Project Authors}
\date{stacks-project:sites-modules.tex:local-snapshot-20261012; GNU FDL 1.2 or later}
\begin{document}
\maketitle
\noindent\textit{Editorial note.} All delivered mathematical text is the author's own, copied verbatim from the pinned Stacks Project chapter file sites-modules.tex; no statement or proof step was rewritten, reordered or removed.\par\medskip
\noindent\textit{Editorial note.} Span recomputation by label pairing (LANE\_BRIEF 10.2): in this chapter the author's proof environment follows the statement environment as a standalone block; the spec covers the complete environment pair (statement span, proof span) as recorded in the pool candidate and re-verified against the pinned source (sm-twins-20261014.json span sanity: all OK).\par\medskip
\noindent\textit{Editorial note.} Occupation scan: combined claim\_key + statement\_anchor + generic-line scan over all registration files (active and withdrawn) found hard occupation only from the 47 pre-existing lane-s11 S11-sites-mod-* cards (same pinned snapshot, delivered) and cross-chapter citation noise, 0 other-lane hits for sites-modules.tex (sm-occupied.json, 2026-10-14); the RESUME-pool dispatch unlocks the remaining face.\par\medskip
\ifdefined\hypertarget\hypertarget{sourcepointer-1}{}\fi
\noindent\textit{Source reference.} Original lemma 0E8J: sources/stacks-chow/tags.txt, source lines 17819--17819. Cross-chapter standard prerequisite cited by the author inside the counted span; it is used as a pointer and its own proof is not part of this item.\par\medskip
\ifdefined\hypertarget\hypertarget{sourcepointer-2}{}\fi
\noindent\textit{Source reference.} Original reference 04VX: sources/stacks-chow/tags.txt, source lines 6340--6340. Cross-chapter standard prerequisite cited by the author inside the counted span; it is used as a pointer and its own proof is not part of this item.\par\medskip
\ifdefined\hypertarget\hypertarget{sourcepointer-3}{}\fi
\noindent\textit{Source reference.} Original reference 00HA: sources/stacks-chow/tags.txt, source lines 636--636. Cross-chapter standard prerequisite cited by the author inside the counted span; it is used as a pointer and its own proof is not part of this item.\par\medskip
\ifdefined\hypertarget\hypertarget{sourcepointer-4}{}\fi
\noindent\textit{Source reference.} Original reference 00DD: sources/stacks-chow/tags.txt, source lines 493--493. Cross-chapter standard prerequisite cited by the author inside the counted span; it is used as a pointer and its own proof is not part of this item.\par\medskip
\begin{lemma}
\label{lemma-flat-eq}
Let $(\mathcal{C}, \mathcal{O})$ be a ringed site. Let $\mathcal{F}$ be an
$\mathcal{O}$-module. The following are equivalent
\begin{enumerate}
\item $\mathcal{F}$ is a flat $\mathcal{O}$-module.
\item Let $U$ be an object of $\mathcal{C}$ and let
$$
\mathcal{O}_U \xrightarrow{(f_1, \ldots, f_n)}
\mathcal{O}_U^{\oplus n} \xrightarrow{(s_1, \ldots, s_n)}
\mathcal{F}|_U
$$
be a complex of $\mathcal{O}_U$-modules. Then there exists a covering
$\{U_i \to U\}$ and for each $i$ a factorization
$$
\mathcal{O}_{U_i}^{\oplus n}
\xrightarrow{B_i}
\mathcal{O}_{U_i}^{\oplus l_i} \xrightarrow{(t_{i1}, \ldots, t_{il_i})}
\mathcal{F}|_{U_i}
$$
of $(s_1, \ldots, s_n)|_{U_i}$ such that
$B_i \circ (f_1, \ldots, f_n)|_{U_i} = 0$.
\item Let $U$ be an object of $\mathcal{C}$ and let
$$
\mathcal{O}_U^{\oplus m} \xrightarrow{A}
\mathcal{O}_U^{\oplus n} \xrightarrow{(s_1, \ldots, s_n)}
\mathcal{F}|_U
$$
be a complex of $\mathcal{O}_U$-modules. Then there exists a covering
$\{U_i \to U\}$ and for each $i$ a factorization
$$
\mathcal{O}_{U_i}^{\oplus n}
\xrightarrow{B_i}
\mathcal{O}_{U_i}^{\oplus l_i} \xrightarrow{(t_{i1}, \ldots, t_{il_i})}
\mathcal{F}|_{U_i}
$$
of $(s_1, \ldots, s_n)|_{U_i}$ such that
$B_i \circ A|_{U_i} = 0$.
\end{enumerate}
\end{lemma}

\begin{proof}
Assume (1). Let $\mathcal{I} \subset \mathcal{O}_U$ be the sheaf of ideals
generated by $f_1, \ldots, f_n$. Then $\sum f_j \otimes s_j$ is
a section of $\mathcal{I} \otimes_{\mathcal{O}_U} \mathcal{F}|_U$
which maps to zero in $\mathcal{F}|_U$. As $\mathcal{F}|_U$ is flat
(Lemma \ref{lemma-restriction-flat}) the map
$\mathcal{I} \otimes_{\mathcal{O}_U} \mathcal{F}|_U \to \mathcal{F}|_U$
is injective. Since $\mathcal{I} \otimes_{\mathcal{O}_U} \mathcal{F}|_U$
is the sheaf associated to the presheaf tensor product, we see
there exists a covering $\{U_i \to U\}$ such
that $\sum f_j|_{U_i} \otimes s_j|_{U_i}$ is zero in
$\mathcal{I}(U_i) \otimes_{\mathcal{O}(U_i)} \mathcal{F}(U_i)$.
Unwinding the definitions using Algebra, Lemma \ref{algebra-lemma-relations}
we find $t_{i1}, \ldots, t_{i l_i} \in \mathcal{F}(U_i)$ and
$a_{ijk} \in \mathcal{O}(U_i)$
such that $\sum_j a_{ijk}f_j|_{U_i} = 0$ and
$s_j|_{U_i} = \sum_k a_{ijk}t_{ik}$.
Thus (2) holds.

\medskip\noindent
Assume (2). Let $U$, $n$, $m$, $A$ and $s_1, \ldots, s_n$ as in (3) be given.
Observe that $A$ has $m$ columns. We will prove the assertion of (3)
is true by induction on $m$. For the base case $m = 0$ we
can take the covering $\{U \to U\}$ and for the factorization we use
$\mathcal{O}_U^n \to \mathcal{O}_U^n \to \mathcal{F}|_U$ where
the first map is the identity and the second is $(s_1, \ldots, s_n)$.
Assume $m > 0$. Let $(f_1, \ldots, f_n)$ be the last column of $A$
and apply (2). This gives new diagrams
$$
\mathcal{O}_{U_i}^{\oplus m} \xrightarrow{B_i \circ A|_{U_i}}
\mathcal{O}_{U_i}^{\oplus l_i} \xrightarrow{(t_{i1}, \ldots, t_{il_i})}
\mathcal{F}|_{U_i}
$$
but the last column of $A_i = B_i \circ A|_{U_i}$ is zero.
Hence we can apply the induction hypothesis to
$U_i$, $l_i$, $m - 1$, the matrix consisting of
the first $m - 1$ columns of $A_i$, and $t_{i1}, \ldots, t_{il_i}$
to get coverings $\{U_{ij} \to U_j\}$ and factorizations
$$
\mathcal{O}_{U_{ij}}^{\oplus l_i}
\xrightarrow{C_{ij}}
\mathcal{O}_{U_{ij}}^{\oplus k_{ij}}
\xrightarrow{(v_{ij1}, \ldots, v_{ij k_{ij}})}
\mathcal{F}|_{U_{ij}}
$$
of $(t_{i1}, \ldots, t_{il_i})|_{U_{ij}}$ such that
$C_i \circ B_i|_{U_{ij}} \circ A|_{U_{ij}} = 0$.
Then $\{U_{ij} \to U\}$ is a covering and we get the
desired factorizations
using $B_{ij} = C_i \circ B_i|_{U_{ij}}$ and
$v_{ija}$. In this way we see that (2) implies (3).

\medskip\noindent
Assume (3). Let $\mathcal{G} \to \mathcal{H}$ be an injective homomorphism
of $\mathcal{O}$-modules. We have to show that
$\mathcal{G} \otimes_\mathcal{O} \mathcal{F} \to
\mathcal{H} \otimes_\mathcal{O} \mathcal{F}$
is injective. Let $U$ be an object of $\mathcal{C}$
and let $s \in (\mathcal{G} \otimes_\mathcal{O} \mathcal{F})(U)$
be a section which maps to zero in
$\mathcal{H} \otimes_\mathcal{O} \mathcal{F}$.
We have to show that $s$ is zero. Since
$\mathcal{G} \otimes_\mathcal{O} \mathcal{F}$
is a sheaf, it suffices to find a covering $\{U_i \to U\}_{i \in I}$
of $\mathcal{C}$ such that $s|_{U_i}$ is zero for all $i \in I$.
Hence we may always replace $U$ by the members of a covering.
In particular, since $\mathcal{G} \otimes_\mathcal{O} \mathcal{F}$
is the sheafification of $\mathcal{G} \otimes_{p, \mathcal{O}} \mathcal{F}$
we may assume that $s$ is the image of $s' \in 
\mathcal{G}(U) \otimes_{\mathcal{O}(U)} \mathcal{F}(U)$.
Arguing similarly for $\mathcal{H} \otimes_\mathcal{O} \mathcal{F}$
we may assume that $s'$ maps to zero in
$\mathcal{H}(U) \otimes_{\mathcal{O}(U)} \mathcal{F}(U)$.
Write $\mathcal{F}(U) = \colim M_\alpha$ as a filtered colimit of finitely
presented $\mathcal{O}(U)$-modules $M_\alpha$
(Algebra, Lemma \ref{algebra-lemma-module-colimit-fp}).
Since tensor product commutes with filtered colimits
(Algebra, Lemma \ref{algebra-lemma-tensor-products-commute-with-limits})
we can choose an $\alpha$ such that $s'$
comes from some $s'' \in \mathcal{G}(U) \otimes_{\mathcal{O}(U)} M_\alpha$
and such that $s''$ maps to zero in
$\mathcal{H}(U) \otimes_{\mathcal{O}(U)} M_\alpha$.
Fix $\alpha$ and $s''$.
Choose a presentation
$$
\mathcal{O}(U)^{\oplus m} \xrightarrow{A} \mathcal{O}(U)^{\oplus n}
\to M_\alpha \to 0
$$
We apply (3) to the corresponding complex of $\mathcal{O}_U$-modules
$$
\mathcal{O}_U^{\oplus m} \xrightarrow{A}
\mathcal{O}_U^{\oplus n} \xrightarrow{(s_1, \ldots, s_n)}
\mathcal{F}|_U
$$
After replacing $U$ by the members of the covering $U_i$
we find that the map
$$
M_\alpha \to \mathcal{F}(U)
$$
factors through a free module $\mathcal{O}(U)^{\oplus l}$ for some $l$.
Since $\mathcal{G}(U) \to \mathcal{H}(U)$ is injective
we conclude that
$$
\mathcal{G}(U) \otimes_{\mathcal{O}(U)} \mathcal{O}(U)^{\oplus l}
\to
\mathcal{H}(U) \otimes_{\mathcal{O}(U)} \mathcal{O}(U)^{\oplus l}
$$
is injective too. Hence as $s''$ maps to zero in the module on
the right, it also maps to zero in the module on the left, i.e.,
$s$ is zero as desired.
\end{proof}

\end{document}
